% Source: https://artofproblemsolving.com/community/c3h1613164_sprint_2017_number_25
% Stable identity: aops-topic-1613164-q25
On an infinitely large grid of squares, a two-by-two region is shaded. Starting in the square above the upper left shaded square, the squares of the grid are numbered with consecutive positive integers, moving in a clockwise direction to the next open square, so that the numbers spiral outwards around the shaded region. The figure shows the first $18$ numbered squares. What integer appears in the $6$th square to the right of the square numbered $5$?

\begin{asy}
usepackage("amsmath");
size(5.5cm);
defaultpen(fontsize(10pt));
defaultpen(fontsize(9pt));
for(int i=-3;i<=5;++i) draw((i,-3)--(i,4),gray(0.7));
for(int j=-3;j<=4;++j) draw((-3,j)--(5,j),gray(0.7));
fill((1,1)--(3,1)--(3,-1)--(1,-1)--cycle,gray(0.55));
draw((1,1)--(3,1)--(3,-1)--(1,-1)--cycle,linewidth(0.9));
int[][] c={{0,1},{1,1},{2,1},{2,0},{2,-1},{2,-2},{1,-2},{0,-2},{-1,-2},{-1,-1},{-1,0},{-1,1},{-1,2},{0,2},{1,2},{2,2},{3,2},{3,1}};
for(int k=0;k<18;++k) label("$"+string(k+1)+"$",(c[k][0]+1.5,c[k][1]+0.5));
label("\mbox{}",(min(currentpicture,user=true)+max(currentpicture,user=true))/2,black); // keeps the text after the figure black when LaTeX typesets the labels inline
\end{asy}
